\chapter{Methodik und Entwicklung}
\label{sec:methodik}

In diesem Kapitel wird die konzeptionelle Entwicklung und technische Umsetzung des vorgestellten Erkennungssystems dieser Arbeit erläutert. Es wird der Weg von der Problemstellung bis hin zur Lösung beschrieben. Dabei wird genauer auf die Designentscheidungen und die Architektur der einzelnen Komponenten der Verarbeitungskette eingegangen.

\section{Gesamtarchitektur und Datenfluss}
Um zuverlässig beurteilen zu können ob eine Katze ein Beutetier im Maul trägt, reicht weder eine reine Objekterkennung noch eine reine Bildsegmentierung aus. Da wir nur einen begrenzten Datensatz an Katzen mit Beutetieren zur Verfügung haben, ist es unerlässlich diese Aufgabe auf mehrere Stufen zu teilen \cite{baumann2020smartcatflap}. Daraus ergibt sich eine Verarbeitungskette aus Objekterkennung, Objektverfolgung, Bildsegmentierung und zeitlicher Zustandsmodellierung. Dabei sieht die finale Verarbeitungskette wie folgt aus:

\begin{enumerate}
	\item Ein Einzelbild der Kamera wird aufgenommen.
	\item Die YOLO Objekterkennung analysiert das Bild.
	\begin{enumerate}
		\item Wenn keine Katze gefunden wurde, endet die Kette hier.
		\item Wird mindestens eine Katze gefunden, so wird jedes Katzengesicht aus dem Bild extrahiert.
	\end{enumerate}
	\item Der ByteTrack-Algorithmus ordnet den Katzengesichtern eine Identität zu oder erstellt eine neue, sollte die Katze frisch ins Bild gekommen sein.
	\item Das YOLO-Segmentierungsmodell untersucht jeden Bildausschnitt einzeln und liefert für jede Katze einen Konfidenzwert, dass sie Beute trägt.
	\item Der Zustandsautomat aggregiert die Ergebnisse der Bildsegmentierung über die Zeit und trifft dann eine finale Entscheidung, in welchem Zustand (z.\,B. \glqq Katze mit Beute\grqq) sich die verfolgten Objekte befindet.
	\begin{enumerate}
		\item Trägt eine Katze Beute im Maul so wird die Klappe geschlossen.
		\item Sind alle Katzen über das gesamte Zeitfenster als sauber erkannt worden, so wird die Klappe geöffnet.
	\end{enumerate}
\end{enumerate}

\section{Vorverarbeitung für Infrarot-Nachtsicht}
Da Katzen ihre Beute nicht nur am Tag bei optimaler Belichtung ins Haus bringen, ist eine zentrale Anforderung an das System das es auch Nachts und bei schlechter Belichtung funktioniert. Das bringt eine weitere Schwierigkeit mit sich: Nachtsichtkameras schalten bei Dunkelheit auf Infrarot um, wodurch Graustufenbilder entstehen.

Damit das neuronale Netz die Farbinformationen nicht erlernt, werden alle Bilder vor dem Training und der Inferenz zu Graustufenbildern konvertiert. Andernfalls würde weder Objekterkennung noch Bildsegmentierung bei Nacht zuverlässig funktionieren.

\section{Training und Anpassung des Objektdetektors (YOLO26n)}
Für die Objekterkennung kommt die Archtitektur YOLO26 in der Ausführung \texttt{yolo26n} zum Einsatz. Sie ist die kleinste und bietet hohe Erkennungsgenauigkeit während es die notwendige Inferenzgeschwindigkeit für Einplatinencomputer liefert.

Standardmodelle von YOLO26 sind nicht spezifisch auf die Erkennung von Katzengesichtern abgestimmt. Daher wurde eine Feinabstimmung für genau dieses Problem durchgeführt \cite{yosinski2014transferable}. Grundlage hierfür ist der Kaggle-Datensatz (Cat Dataset von Crawford \cite{kaggle2024cat}), da dieser bereits Annotationen für das Gesicht enthält. Der Datensatz beinhaltet die Koordinaten für Augen, Ohren und Mund. Um diese für YOLO nutzbar zu machen, müssen sie in ein Begrenzungsrahmen-Format umgewandelt werden. Dafür wurde ein Algorithmus entwickelt, der die Breite des Rahmens durch die Ohren und die Höhe durch die Ohren und den Mund vorgibt.

Die Feinabstimmung des YOLO26n-Modells erfolgte über 100 Epochen und einer Bildauflösung von 640x640 Pixeln und der Standard YOLO-Augmentation \cite{ultralytics_augmentation}. Um eine Überanpassung des Modells zu vermeiden, wird das Training abgebrochen, sollte in den letzten zehn Epochen keine Verbesserung mehr eintreten. \cite{goodfellow2016deep}.

\section{Automatisierte Generierung des Segmentierungs-Datensatzes}
Für die Segmentierung, ob die Katze Beute im Maul trägt, gab es leider keinen brauchbaren Datensatz. Für das Training werden zwei Arten von Bildern benötigt, Katzen mit und Katzen ohne Beute. Die Katzen ohne Beute können wir wieder aus dem Kaggle-Datensatz (Cat Dataset von Crawford \cite{kaggle2024cat}) verwenden. Katzen mit Beute auf der anderen Seite sind extrem schwierig zu bekommen, entweder sie sind qualitativ minderwertig oder Urheberrechtlich geschützt. 

Die Bilder von Katzen mit Beute, lies ich mir daher von Gemini \cite{geminiteam2023gemini} generieren. Mithilfe der offiziellen Gemini Application Programming Interface (API), entwickelte ich einen Algorithmus, der unter Berücksichtigung des Benutzungslimits automatisch Bilder von Katzen mit Beute in unterschiedlichen Posen und Entfernungen generiert.

Für die Vorbereitung des Datensatzes zur Feinabstimmung des YOLO26n Segmentierungsmodells, wird der zuvor trainierte Katzengesichts-Detektor verwendet. Er erkennt die Katzengesichter in den Bildern ohne Beute und schneidet sie genauso zu wie später bei der Inferenz. Diese Ausschnitte werden als Negativbeispiele ohne Annotation in den Datensatz aufgenommen. Die Bilder mit Beute wurden mithilfe der KI-gestützten Annotation von Roboflow \cite{roboflow2025} mit Segmentierungsmasken der Beutetiere versehen und anschließend händisch überprüft. 

Das Segmentierungsmodell hat somit immer die gleiche Eingangslage: Es bekommt stets den erweiterten Ausschnitt um das Katzengesicht ohne viel Hintergrund. Das bildet die perfekte Trainingsgrundlage, da die Segmentierung in der späteren Inferenz genau die gleichen Eingangsbilder erhält.

\section{Bildausschnitt-Extraktion}
Da der Eingang des Segmentierungsmodells der Ausgang des Detektors ist, müssen beide die erforderlichen Informationen erhalten. Während der Detektor am besten funktioniert, wenn er nur auf das Katzengesicht trainiert wird, benötigt er auch die Information, was sich unter dem Maul des Katzengesichts befindet.

Daher wird der Begrenzungsrahmen des Detektors verändert bevor ihn das Segmentierungsmodell bekommt. Das Katzengesicht wird nach untenhin um die doppelte Gesichtshöhe und seitlich um jeweils 80\,\% der Gesichtsbreite erweitert um sicher zu gehen das die Beute im Bild ist.

Der Detektor erwartet ein 3-Kanal-Bild, daher wird das Graustufenbild dreifach dupliziert (\texttt{cv2.cvtColor(gray, cv2.COLOR\_GRAY2BGR)}). Das Segmentierungsmodell wurde dagegen direkt mit einkanaligen Graustufenbildern trainiert und erhält den Ausschnitt als einkanaliges Bild. Zuvor wird der Ausschnitt mit schwarzen Rändern zu einem Quadrat aufgefüllt.

\section{Training des Segmentierungsmodells (YOLO-Segmentierung)}
Als Segmentierungsmodell kommt wie beim Detektor ein YOLO-Modell zum Einsatz. Dadurch das wir uns bereits in der Ultralytics Umgebung befinden und moderne YOLO-Versionen hochoptimierte Netzwerke für Segmentierungsaufgaben bereitstellen, war diese Entscheidung naheliegend.

Das Modell bekommt ein 640x640 Pixel großes Bild und liefert für jedes erkannte Beutetier eine Segmentierungsmaske und einen Konfidenzwert. Dadurch dass der Datensatz des Modells relativ klein ist, wird hier stark auf Datenvervielfältigung gesetzt (zufällige Rotation des Bildes um bis zu 10°, horizontale Spiegelung, Variation von Helligkeit und Sättigung). Die Mosaik-Augmentation, bei der vier Bilder zu einem zusammengesetzt werden, ist deaktiviert, damit kleine Beutetiere nicht verloren gehen. Durch diese Augmentationsschritte wird das Modell robuster gegen unbekannte Posen der Katzen. 

\section{Objektverfolgung und Entwicklung des Zustandsmodells}
Wenn sich eine Katze der Klappe nähert, ändert sich die Qualität der Bilder durch Bewegungsunschärfe oder ungünstige Winkel stark. Das hat zur Folge, dass die Ausgaben des Segmentierungsmodells schwanken. Wenn z.b. die Beute kurzzeitig verdeckt ist, würde das zur Folge haben das die Katze ins Haus darf. 

Damit dieses Problem eingedämpft wird, wurde die Objekterkennung um den ByteTrack-Algorithmus erweitert. Dadurch bekommt eine Katze beim ersten auftreten im Bild eine eindeutige Objekt-ID, diese ID bleibt in den darauffolgenden Bildern für diese Katze gleich. Basierend auf der Eindeutigkeit jeder Katze konnte nun ein Zustandsautomat entwickelt werden. Diese Klasse \texttt{StateMachine} verwaltet für jede erkannte Katze (\texttt{track\_id}) ein \texttt{TrackState}-Objekt. Das Modell kennt vier Zustände:

\begin{itemize}
    \item \texttt{NO\_CAT}: Es befindet sich kein relevantes Objekt im Sichtfeld.
    \item \texttt{GATHERING\_HISTORY}: Eine Katze wurde erkannt, es liegen aber noch nicht genügend Einzelbilder für eine Entscheidung vor.
    \item \texttt{CAT\_NO\_PREY}: Eine Katze wurde erkannt, trägt jedoch keine Beute.
    \item \texttt{CAT\_WITH\_PREY}: Eine Katze wurde erkannt und trägt mit hinreichender Sicherheit eine Beute.
\end{itemize}

Damit nun die Systementscheidung nicht durch das aktuellste Einzelbild getroffen wird, speichert das \texttt{TrackState}-Objekt eine Historie der Konfidenzwerte des Segmentierungsmodells der letzten 20 Bilder. Der aktuelle Zustand wird bei jedem Bild neu berechnet, dabei wird einfach der gleitende Durchschnitt der Historie berechnet. Wird der Schwellenwert von 0,3 erreicht oder überschritten, so wird die Katze als \texttt{CAT\_WITH\_PREY} markiert. Fällt der Wert darunter wird sie als \texttt{CAT\_NO\_PREY} markiert, sofern bereits 20 Werte gesammelt wurden.

Der Schwellenwert von 0,3 wurde anhand der Parameterreihe in Abschnitt~\ref{sec:ablation_sm} gewählt. Da der Durchschnitt über 20 Einzelbilder einzelne Ausreißer bereits stark glättet, kann der Schwellenwert niedrig angesetzt werden, ohne dass die Katze zu Unrecht draußen bleiben muss. Dieser Wert ist aber nicht in Stein gemeißelt. Wenn es für den Endanwender wichtiger wäre, dass die Katze im Zweifel ins Haus darf, würde man den Wert höher ansetzen, beispielsweise bei 0,5.

Aus den Zuständen aller Katzen ergibt sich die Entscheidung über die Klappe. Trägt mindestens eine Katze Beute, wird die Klappe sofort gesperrt. Freigegeben wird sie erst, wenn alle Katzen im Bild eine vollständige Historie besitzen und als sauber gelten. Dabei werden nur Katzen berücksichtigt, die im aktuellen Einzelbild sichtbar sind. Vergibt der Tracker einer Katze eine neue Identität, blockiert der zurückgebliebene, kurzlebige Track die Freigabe somit nicht. Eine Katze mit Beute hält die Sperre hingegen aufrecht, bis sie aus dem Speicher entfernt wird. Damit die Klappe nicht ständig auf- und zugeht, wird eine erneute Freigabe frühestens fünf Sekunden nach dem letzten Wechsel erlaubt.

Um Speicherlecks in der Zustandsmaschine zu verhindern, werden Katzen die das Bildfeld verlassen haben oder die länger als 30 Bilder nicht mehr gefunden werden (\texttt{max\_missed\_frames}) aus dem Speicher entfernt.
